\textbf{Proof.}
Ass:exposure-positivity makes every occurrence of \(\Pi^{-1}\) well defined, and ass:frechet-covariance-domain supplies the feasible endpoints \(c_L,c_U\).  Define
\[
V_+=\inf_{[B]_+\in\mathcal B_Z^+}w^+(B),
\qquad
V_Z=\inf_{C\in\mathcal B_Z}w(C).                                  \tag{1}
\]
These are infima; no minimizer is presumed.

For every \(C\in\mathcal B_Z\), lem:zhao-rank-two-stratum-identification gives
\[
w^+(\iota(C))=w(C).
\]
Since \(\iota(C)\in\mathcal B_Z^+\), it follows that
\[
V_+\le V_Z.                                                        \tag{2}
\]

For the reverse inequality, fix any \([B]_+\in\mathcal B_Z^+\).  If \(\operatorname{rank}(S_B)=2\), the same lemma supplies a unique \(C\in\mathcal B_Z\) with \([B]_+=\iota(C)\), and therefore
\[
V_Z\le w(C)=w^+(B).                                                \tag{3}
\]
If \(\operatorname{rank}(S_B)\le1\), lem:rank-stratified-projector-reduction supplies \(C_j\in\mathcal B_Z\) with
\[
\limsup_{j\to\infty}w(C_j)\le w^+(B).
\]
Because \(V_Z\le w(C_j)\) for every \(j\),
\[
V_Z\le\liminf_{j\to\infty}w(C_j)
\le\limsup_{j\to\infty}w(C_j)
\le w^+(B).                                                        \tag{4}
\]
Thus \(V_Z\le w^+(B)\) for every completion class.  Taking the infimum in (3)--(4) gives \(V_Z\le V_+\), which together with (2) proves
\[
v_{\mathrm{mm}}^+=V_+=V_Z.                                       \tag{5}
\]
Prop:kappa-value-exhaustion then yields
\[
v_{\mathrm{mm}}^+
=\inf_{C\in\mathcal B_Z}w(C)
=\inf_{\substack{\kappa>0:\,\mathcal B_\kappa\ne\varnothing}}
\min_{C\in\mathcal B_\kappa}w(C).                               \tag{6}
\]
The final index set is nonempty because \(\mathcal B_Z\ne\varnothing\) and prop:kappa-exhaustion places every ordinary class in some threshold class.  For each nonempty fixed threshold, compactness from prop:kappa-exhaustion and continuity of the ordinary-inverse loss under the eigenvalue bound \(S_B\succeq n\kappa I_2\) justify the minimum in (6).

We next identify the equality set.  Suppose \(\iota(C)\in\mathcal A_{\mathrm{mm}}^+\).  By (5) and rank-two equality,
\[
w(C)=w^+(\iota(C))=v_{\mathrm{mm}}^+=V_Z,                          \tag{7}
\]
so \(C\) attains the ordinary-domain infimum.  Prop:kappa-exhaustion supplies \(\kappa_C>0\) with \(C\in\mathcal B_{\kappa_C}\).  Equation (6) and (7) give
\[
v_{\mathrm{mm}}^+
\le\min_{A\in\mathcal B_{\kappa_C}}w(A)
\le w(C)=v_{\mathrm{mm}}^+.
\]
Hence \(C\in\mathcal A_{\mathrm{mm}}(\kappa_C)\) and the threshold minimum equals \(v_{\mathrm{mm}}^+\).

Conversely, if \(C\in\mathcal A_{\mathrm{mm}}(\kappa)\) for a nonempty threshold class satisfying
\[
\min_{A\in\mathcal B_\kappa}w(A)=v_{\mathrm{mm}}^+,
\]
then rank-two equality gives \(w^+(\iota(C))=w(C)=v_{\mathrm{mm}}^+\).  Therefore \(\iota(C)\in\mathcal A_{\mathrm{mm}}^+\).  The two implications prove
\[
\mathcal A_{\mathrm{mm}}^+\cap\iota(\mathcal B_Z)
=\iota\!\left(
\bigcup_{\substack{\kappa>0:\,
\mathcal B_\kappa\ne\varnothing,\ 
\min_{B\in\mathcal B_\kappa}w(B)=v_{\mathrm{mm}}^+}}
\mathcal A_{\mathrm{mm}}(\kappa)\right).                          \tag{8}
\]
They also prove the two conditional optimizer-transfer assertions in the theorem.

Finally, lem:rank-stratified-projector-reduction partitions the completion into ranks zero, one, and two, while lem:zhao-rank-two-stratum-identification identifies rank two exactly with \(\iota(\mathcal B_Z)\).  Removing that stratum from the equality set gives
\[
\mathcal A_{\mathrm{mm}}^+\setminus\iota(\mathcal B_Z)
=\{[B]_+\in\mathcal B_Z^+:\operatorname{rank}(S_B)\le1,
\ w^+(B)=v_{\mathrm{mm}}^+\}.                                    \tag{9}
\]
For a singular equality point, the projector lemma produces an ordinary minimizing sequence, not an ordinary minimizer.  Accordingly, (1)--(9) establish neither unconditional completion attainment nor ordinary attainment from a singular equality point, and they use no cross-rank continuity. \(\square\)

\textbf{Sanity check.}
If an ordinary minimizer exists, its image satisfies the equality defining \(\mathcal A_{\mathrm{mm}}^+\) by (5).  If no ordinary minimizer exists, the infimum equality (5) remains valid, while the completion equality set may consistently be empty.